%HEADER%
\documentclass[11pt,a4paper]{article}
% Le préambule commun à toutes les transcriptions du fonds Grothendieck.
%
% Il tient en une idée : **l'incertitude se compose, elle ne se résout pas.**
% Une transcription qui lisse un mot douteux a détruit la seule chose pour
% laquelle on la faisait — un lecteur doit pouvoir savoir, à chaque endroit,
% ce qui était sur le feuillet et ce qui a été ajouté. Les six macros qui
% suivent sont donc l'appareil critique, et non de la décoration.
%
% Les mêmes commandes sont reconnues par `scripts/render.mjs`, qui produit la
% vue de lecture du site. Une macro ajoutée ici doit l'être aussi là-bas,
% faute de quoi le rendu échouera bruyamment — ce qui est voulu : un rendu
% silencieusement faux est pire qu'un rendu absent.

\ProvidesPackage{grothendieck}[2026/08/08 Transcriptions du fonds Grothendieck]

\usepackage{amsmath,amssymb,amsthm}
\usepackage{mathrsfs}
\usepackage{stmaryrd}
% Les diagrammes commutatifs sont partout dans ce fonds : c'est la langue
% même de la géométrie algébrique de Grothendieck, et une transcription qui
% les rendrait en prose ne transcrirait rien.
\usepackage{tikz-cd}
% Français par défaut — c'est la langue des feuillets — mais l'anglais est
% chargé aussi : la traduction et le résumé sont des documents anglais, et la
% typographie française (espaces avant les deux-points) y serait une faute.
% Ces documents basculent par \selectlanguage{english} en tête de corps.
\usepackage[english,french]{babel}
\usepackage{fontspec}
\usepackage{geometry}
\geometry{a4paper,margin=25mm,marginparwidth=28mm}
\usepackage{marginnote}
\usepackage{xcolor}
% ulem, not soul. `soul`'s \st and \ul are built on a character-by-character
% scan that XeLaTeX's font machinery breaks: the document fails with an
% undefined control sequence rather than degrading. ulem does the same two jobs
% and is Unicode-engine safe. `normalem` stops it hijacking \emph, which the
% transcriptions use for Grothendieck's own emphasis.
\usepackage[normalem]{ulem}
\usepackage{hyperref}
\hypersetup{colorlinks=true,linkcolor=black,urlcolor=black,pdfborder={0 0 0}}

% --- Métadonnées du lot -----------------------------------------------------
% Reprises telles quelles par le script de rendu, qui les affiche en tête de
% la vue de lecture. Elles fixent ce que le fichier prétend couvrir : sans
% elles, une transcription flotte sans point d'attache dans un fonds de seize
% mille feuillets.

\newcommand{\folder}[1]{\gdef\grothFolder{#1}}
\newcommand{\batch}[1]{\gdef\grothBatch{#1}}
\newcommand{\leaves}[2]{\gdef\grothFirst{#1}\gdef\grothLast{#2}}
\newcommand{\foldertitle}[1]{\gdef\grothFoldertitle{#1}}
\gdef\grothFolder{?}\gdef\grothBatch{?}\gdef\grothFirst{?}\gdef\grothLast{?}\gdef\grothFoldertitle{}

% --- Le résumé --------------------------------------------------------------
%
% Il ouvre la lecture modernisée, sous le titre « Résumé », et il est composé
% en petit. Ce n'est pas un ornement : c'est le seul passage du document qui
% ne soit pas des mathématiques, et un lecteur doit pouvoir le reconnaître
% avant de l'avoir lu — puis le sauter s'il connaît déjà le sujet, ou s'y
% arrêter s'il ne le connaît pas. Le filet qui le suit marque la reprise du
% corps.

\newenvironment{resume}
  {\small\setlength{\parskip}{0.45em}}
  {\par\medskip\hrule\medskip}

% --- Mots-clés ---------------------------------------------------------------
%
% En anglais, en clôture du résumé de la lecture modernisée : le vocabulaire
% moderne sous lequel chercher ce que les feuillets construisent. Le manifeste
% du site (`npm run manifest`) les extrait pour étiqueter la cote entière —
% cette ligne est la source unique des tags, il n'y a pas de fichier à côté.
\newcommand{\keywords}[1]{\par\smallskip\noindent
  {\footnotesize\textsc{Keywords} --- \textit{#1}}\par}

% --- L'appareil critique ----------------------------------------------------

% Le numéro de feuillet, en marge. C'est l'ancre qui permet de régler un
% désaccord sur une formule en regardant *un* feuillet plutôt qu'une chemise
% de six cents. Le site s'en sert aussi pour tourner les pages du fac-similé
% quand on fait défiler la transcription.
\newcommand{\leaf}[1]{%
  \marginnote{\footnotesize\color{gray!70}\textsf{#1}}%
  \label{leaf:#1}\ignorespaces}

% Illisible : on ne devine pas. Un mot inventé qui se lit comme les autres est
% le pire résultat possible.
\newcommand{\ill}{\textcolor{red!60!black}{[\,\ldots\,]}}

% Lecture douteuse : le mot est donné, et signalé comme incertain.
\newcommand{\uncertain}[1]{\uline{#1}}

% Ajout de l'éditeur — une lettre manquante, un mot sous-entendu.
\newcommand{\add}[1]{\textcolor{blue!50!black}{[#1]}}

% Biffé par Grothendieck lui-même. Ce qu'il a rayé fait partie du document :
% une rature dit souvent où la pensée a changé de direction.
\newcommand{\struck}[1]{\sout{#1}}

% Note du transcripteur, sur l'état matériel du feuillet.
\newcommand{\note}[1]{\par\smallskip{\small\color{gray!60!black}\textit{[#1]}}\par\smallskip}

% Note marginale de Grothendieck, distincte du corps du texte.
\newcommand{\marginal}[1]{\par\smallskip\noindent\hspace{1em}%
  {\small\color{gray!60!black}#1}\par\smallskip}

% --- Titre ------------------------------------------------------------------

\AtBeginDocument{%
  \begin{center}
    {\large\bfseries Fonds Alexandre Grothendieck --- cote n\textsuperscript{o}~\grothFolder}\\[2pt]
    {\itshape\grothFoldertitle}\\[4pt]
    {\small feuillets \grothFirst--\grothLast\ (lot \grothBatch)}\\[2pt]
    {\footnotesize\color{gray!60!black}Transcription établie d'après le fac-similé numérisé
     par l'Université de Montpellier.\\
     Les lectures incertaines sont soulignées, les illisibles notées [\,\ldots\,],
     les ajouts entre crochets.}
  \end{center}
  \vspace{1em}}

\endinput


\folder{140-1}
\batch{1}
\pages{1}{20}
\dating{[à partir de 1978-à partir de 1982]}
\watermark{Édition de démonstration}
\foldertitle{"Longue Marche" [brouillon] : notes manuscrites (s.d.)}

\begin{document}

\section{« Longue Marche »}

\page{1}

\note{le titre est de sa main, entre guillemets, dans un cadre ouvert en haut à
droite de la page 1 ; c'est le seul titre du lot. Page au crayon}

\[
-v d^2 + w cd + c^2
\]

\note{écrit sous le titre, isolé ; c'est le déterminant calculé plus bas pour
$v_1 = 1$}

\struck{$\underline{u}$} $\alpha = \underline{u}\, a^{-1}$

$\underline{u} = \alpha\, \underline{a}$
\qquad
$\begin{pmatrix} a & b \\ c & d \end{pmatrix}$
\qquad
$a = \alpha^{-1} u$

\[
\det(d\rho + c) = d^2 \det\rho + cd\, \mathrm{Tr}\,\rho + c^2
\]

\note{dans $d\rho + c$, le $c$ est surchargé sur une autre lettre}

\[
\rho = \begin{pmatrix} 0 & v_1 \\ v & w \end{pmatrix}
\qquad
d_0 \rho + c_0 = \begin{pmatrix} c_0 & d_0 v_1 \\ d_0 v & d_0 w + c_0 \end{pmatrix}
\]

\[
\begin{pmatrix} a & b \\ c & d \end{pmatrix}
\begin{pmatrix} c & d \\ dv & dw + \uncertain{c} \end{pmatrix}
= \begin{pmatrix} ac + bdv & \\ & \end{pmatrix}
\]

\note{le produit s'arrête au premier coefficient}

\[
\underline{a} = \alpha^{-1} u = \frac{1}{\delta_\alpha}
\begin{pmatrix} d & -b \\ -c & a \end{pmatrix}
\begin{pmatrix} \mu & \lambda \\ 0 & 1 \end{pmatrix}
= \frac{1}{\delta_\alpha}
\begin{pmatrix} d\mu & d\lambda - b \\ -c\mu & -c\lambda + a \end{pmatrix}
= c_0 \rho + d_0
\]

\note{les deux coefficients $\mu$ et $\lambda$ de la seconde matrice sont
récrits sur d'autres lettres (le second sur un « $B_1$ » ?) ; le $d\mu$ du
produit est écrit par-dessus une rature, et le $c_0$ final sur une autre
lettre}

\[
c_0 = d\mu \struck{/} \delta_\alpha = \frac{-dv}{\uncertain{-v^2} + wcd + c^2}
\]

\struck{$d_0 = \frac{1}{\delta_\alpha}(d\lambda - b) = \frac{1}{\delta_\alpha}\bigl(-v bd^2 + v_1 acd + wbcd\bigr.$}

$d_0 \struck{v} = -\dfrac{c\mu}{\delta_\alpha}$
\qquad
$d_0 = -$

\note{la ligne s'arrête sur le signe}

\[
c_0 = \frac{d\mu}{\delta_\alpha} \qquad d_0 = \frac{-c\mu}{\delta_\alpha v}
\]

\[
\boxed{\underline{a} = \frac{\mu}{\delta_\alpha}\Bigl(-\frac{c}{v}\rho + d\Bigr)}
\]

\note{devant le cadre, un signe raturé souligné ; à gauche, biffé :
« $\det \underline{a} = \det$ » sur une fraction de dénominateur $\delta_\alpha$
illisible}

$\det a \overset{?}{=} \delta_\alpha^{-1} \mu$

\[
\det \underline{a} = \Bigl(\frac{\mu}{\delta_\alpha}\Bigr)^2
\Bigl[\frac{c^2}{v^2}(-v v_1) + \frac{dc}{v} w + d^2\Bigr] \ldots
\]

\note{le signe entre les deux premiers termes du crochet est un $+$ repassé en
tache}

\section{Extensions centrales de $\Pi^{D}_n$}

\note{titre de l'édition. Les pages 3, 5, 7 et 11 portent ses numéros ① à
④ (entourés), la page 13 « 4 bis », les pages 15, 17 et 19 les numéros 5, ⑥
et ⑦ : c'est une suite numérotée de sa main. Les pages 9, 14, 16, 18 et 20
ne portent pas de numéro}

\page{3}\note{p. ① de l'auteur}

\[
\widetilde{\Pi}{}^{\circ D}_{n} \struck{\overset{\text{déf}}{=}}
\{\tilde\rho, \tilde\sigma \mid \tilde\rho^{\,n} = \tilde\sigma^{2}\}
\qquad
\boxed{\tilde\rho^{\,n} = \tilde\sigma^{2} = \tilde\omega_0} \ \text{élément central}
\]

\[
1 \to \mathbb{Z} \to \widetilde{\Pi^{\circ D}_{n}} \to \Pi^{D}_{n} \to 1
\qquad \text{ext. centrale}
\]

\note{le petit $\circ$ en exposant de $\Pi$, ici et à la ligne précédente, est
peu marqué ; sa lecture est incertaine}

Posons
\[
\boxed{\tilde\rho_0 = \tilde\sigma\tilde\rho\tilde\sigma^{-1}\tilde\rho}
= \tilde\sigma(\tilde\rho)\,\tilde\rho
= (\tilde\sigma\tilde\rho)^2\,\tilde\omega_0^{-1}
\qquad
\boxed{\tilde\rho_i = \tilde\rho^{\,i}(\tilde\rho_0)}
\qquad \tilde\rho_i \text{ relève } \rho_i
\]

\note{$\tilde\sigma(\tilde\rho)$, $\tilde\rho^{\,i}(\tilde\rho_0)$ : l'action
par conjugaison. L'exposant de $\tilde\rho$ dans la définition de
$\tilde\rho_i$ est une tache, lue $i$}

On a
\begin{align*}
\tilde\rho_{n-1}\tilde\rho_{n-2}\cdots\tilde\rho_0
&= \underbrace{\tilde\rho^{\,n-1}\tilde\rho_0\tilde\rho^{-(n-1)}}\,
\underbrace{\tilde\rho^{\,n-2}\tilde\rho_0\tilde\rho^{-(n-2)}} \cdots
(\tilde\rho\tilde\rho_0\tilde\rho^{-1})\,\tilde\rho_0 \\
&= \tilde\rho^{\,n-1}\tilde\rho_0\tilde\rho^{-1}\tilde\rho_0\tilde\rho^{-1} \cdots
\tilde\rho^{-1}\tilde\rho_0\tilde\rho^{-1}\tilde\rho_0
\qquad (n \text{ facteurs } \tilde\rho_0) \\
&= \tilde\rho^{\,n}\,\underbrace{(\tilde\rho^{-1}\tilde\rho_0)\cdots(\tilde\rho^{-1}\tilde\rho_0)}_{n \text{ facteurs}}
= \underbrace{\tilde\rho^{\,n}}_{\tilde\omega_0}(\tilde\rho^{-1}\tilde\rho_0)^n
\end{align*}

Or
\[
\tilde\rho^{-1}\tilde\rho_0 = \tilde\rho^{-1}\tilde\sigma\tilde\rho\tilde\sigma^{-1}\tilde\rho
= (\tilde\rho^{-1}\tilde\sigma)(\tilde\rho)
\]
donc
\[
(\tilde\rho^{-1}\tilde\rho_0)^n = (\tilde\rho^{-1}\tilde\sigma)(\tilde\rho^{\,n})
= (\tilde\rho^{-1}\tilde\sigma)(\tilde\omega_0) = \tilde\omega_0
\]
d'où
\[
\boxed{\tilde\rho_{n-1}\tilde\rho_{n-2}\cdots\tilde\rho_0 = \tilde\omega_0^{2}}
\]

Soit $\widetilde{\Pi}_n$ \struck{le sous-groupe de $\widetilde{\Pi^{D}_{n}}$
engendré par les $\tilde\rho_i$, il est invariant par $\rho$, et aussi par
$\sigma$} l'image inverse de $\Pi_n \subset \Pi^{D}_n$ dans
$\widetilde{\Pi^{D}_{n}}$

\note{la correction « l'image inverse de $\Pi_n \subset \Pi^D_n$ dans
$\widetilde{\Pi^D_n}$ » est écrite au-dessus des trois lignes biffées ; à
gauche, une amorce raturée : « \struck{$\widetilde{\Pi}$ \ill{}
$\tilde\rho_{n-1}\tilde\rho_{n-2}$} »}

\[
1 \to \mathbb{Z} \to \widetilde{\Pi}_n \to \Pi_n \to 1
\]
\[
\widetilde{\Pi}_n = \{\tilde\rho_0, \ldots, \tilde\rho_{n-1}, \tilde\omega_0 \mid
\tilde\omega_0 \text{ central},\ \tilde\rho_{n-1}\tilde\rho_{n-2}\cdots\tilde\rho_0 = \tilde\omega_0^{2}\}
\]

Soit
\[
\widetilde{\Pi}{}^{\circ}_n = \{\tilde\rho_0, \ldots, \tilde\rho_{n-1} \mid
\tilde\rho_{n-1}\tilde\rho_{n-2}\cdots\tilde\rho_0 \text{ central i.e. }
\tilde\rho_{n-1}\tilde\rho_{n-2}\cdots\tilde\rho_0
= \tilde\rho_{n-2}\tilde\rho_{n-3}\cdots\tilde\rho_0\tilde\rho_{n-1} = \cdots\}
\]

Alors
\[
1 \to \mathbb{Z} \to \widetilde{\Pi}{}^{\circ}_n \to \Pi_n \to 1
\]
$\mathbb{Z}$ engendré par $\varpi = \tilde\rho_{n-1}\tilde\rho_{n-2}\cdots\tilde\rho_0$
\[
\widetilde{\Pi}_n \simeq \widetilde{\Pi}{}^{\circ}_n \wedge_{\mathbb{Z}\varpi} \mathbb{Z}\omega
\qquad \varpi_0 \longmapsto \omega_0^{2}
\]

\note{au-dessus de $\Pi_n$ dans la suite exacte, un mot court, lu
\uncertain{can} ; l'exposant du premier $\widetilde{\Pi}_n$ de la dernière ligne
est une tache}

\page{5}\note{p. ② de l'auteur}

$\widetilde{\Pi}{}^{\circ}_n$ invariant dans $\widetilde{\Pi}_n$ (ss-groupe
d'indice 2)\marginal{ss-gr. engendré par les $\tilde\rho_i$} est-il invariant
dans $\widetilde{\Pi}{}^{D}_n$ ? Invariant par $\tilde\rho$ c'est clair --- donc
invariant dans $\widetilde{\Pi}{}^{D+}_n$. Invariant par $\tilde\sigma$ ?

\note{la note marginale, à gauche, est reliée par un arc à
$\widetilde{\Pi}{}^{\circ}_n$}

dans $\Pi^D_n$ :
\[
\sigma^{-1}(\rho_0) = \rho_1,\quad \sigma^{-1}(\rho_1) = \rho_0,\quad
\sigma^{-1}(\rho_2) = \rho_0(\rho_{n-1}),\ \ldots\
\sigma^{-1}(\rho_i) = \rho_n\cdots\rho_{n+2-i}(\rho_{n+1-i})\ \ldots
\]

dans $\widetilde{\Pi}{}^{D}_n$ :
\[
\tilde\sigma^{-1}(\tilde\rho_0) = \tilde\rho_1\tilde z_0,\quad
\tilde\sigma^{-1}(\tilde\rho_1) = \tilde\rho_0\tilde z_1\ \cdots\
\tilde\sigma^{-1}(\tilde\rho_i) = (\tilde\rho_n\cdots\tilde\rho_{n+2-i})(\tilde\rho_{n+1-i})\,\tilde z_i
\]
\[
z_{n-1}z_{n-2}\cdots z_0 = 1
\]

dans $(\widetilde{\Pi}{}^{D}_n)_{\mathrm{ab}}$ :
\[
\tilde\sigma^{-1}(\tilde\rho_0) = \rho_1 + \struck{\ill{}}\, z_0,\quad
\tilde\sigma^{-1}(\rho_1) = \rho_0 + z_1,\ \ldots\
\tilde\sigma^{-1}(\rho_i) = \rho_{n+1-i} + z_i
\]

\note{le premier exposant $-1$ de la ligne de $\Pi^D_n$ est récrit sur une
rature ; dans la dernière formule de cette ligne, les indices sont surchargés
et se lisent mal}

$z_i = \tilde\omega_0^{\,\nu_i}$ la condition d'invariance de
$\widetilde{\Pi}{}^{\circ}_n$ est que les $\nu_i \equiv 0\ (2)$. On va \ill{}
de \struck{\ill{}} calculs au moins $z_0$, $z_1$.

\[
\tilde\sigma(\tilde\rho_0) = \tilde\sigma^{-1}(\tilde\rho_0)
= \tilde\rho\,\tilde\sigma^{-1}\tilde\rho\,\tilde\sigma
\qquad
\bigl(\tilde\rho_0 = \tilde\sigma\tilde\rho\tilde\sigma^{-1}\tilde\rho\bigr)
\]
\[
\tilde\rho_1 = \tilde\rho\tilde\sigma\tilde\rho\tilde\sigma^{-1}
\]
\[
\tilde\rho_1^{-1}\,\tilde\sigma^{-1}(\tilde\rho_0)
= \tilde\sigma\tilde\rho^{-1}\tilde\sigma^{-1}\tilde\rho^{-1}\tilde\rho\,\tilde\sigma^{-1}\tilde\rho\tilde\sigma
= \tilde\sigma\underbrace{\tilde\rho^{-n}}_{\tilde\omega_0^{-1}}\tilde\sigma
= \struck{z_0}\ \tilde\omega_0^{-1}\underbrace{\tilde\sigma^{2}}_{\tilde\omega_0} = 1
\]

\note{dans la première ligne, l'exposant $-1$ du second $\tilde\sigma$ est
ajouté en appuyant. Dans le calcul de $\tilde\rho_1^{-1}\tilde\sigma^{-1}(\tilde\rho_0)$,
des accolades sous le membre de droite regroupent $\tilde\rho^{-1}\tilde\rho$,
puis $\tilde\sigma^{-1}\cdots\tilde\sigma^{-1}$ en
« $\tilde\sigma^{-2} = \tilde\rho^{-n}$ », le tout noté $\tilde\rho^{-n}$ ; sous
le résultat, un signe raturé}

\struck{$\tilde\sigma(\tilde\rho_1)$}
\[
\tilde\sigma^{-1}(\tilde\rho_1) = \tilde\sigma^{-1}(\tilde\rho_1)
= \tilde\sigma^{-1}\tilde\rho\tilde\sigma\tilde\rho
\]
\[
\tilde\rho_0^{-1}\,\tilde\sigma^{-1}(\tilde\rho_1)
= \tilde\rho^{-1}\tilde\sigma\,\tilde\rho^{-1}\tilde\sigma^{-1}\tilde\sigma^{-1}\tilde\rho\tilde\sigma\tilde\rho = 1
\]

\note{sous ce produit, trois accolades emboîtées, étiquetées $\omega_0^{-2}$,
\uncertain{$\omega_0^{-2}$} et $1$ ; le premier facteur est récrit sur une
rature, et la lecture du deuxième ($\tilde\sigma$ ou $\tilde\sigma^{-1}$) est
incertaine}

\[
\struck{\ill{}}\ \tilde\rho_i = \tilde\rho^{\,i}\tilde\sigma\tilde\rho\tilde\sigma^{-1}\tilde\rho^{\,1-i}
\]
\[
\tilde\sigma(\tilde\rho_i) = \tilde\sigma\tilde\rho^{\,i}\tilde\sigma\tilde\rho\tilde\sigma^{-1}\tilde\rho^{\,1-i}\tilde\sigma^{-1}
\]
\[
\tilde\rho_n\cdots\tilde\rho_{n+2-i}\,\tilde\rho_{n+1-i}\,\tilde\rho_{n+2-i}^{-1}
\]

\note{la page s'arrête sur cette dernière expression}

\page{7}\note{p. ③ de l'auteur ; feuillet en largeur}

\begin{tikzcd}[column sep=small, row sep=small, nodes={font=\scriptsize}]
 & & \widetilde{\Pi}{}^{\circ}_{0,3} \arrow[dr, "\mathbb{Z}/2\mathbb{Z}"] & & \widetilde{\mathrm{SL}}(2,\mathbb{Z}) \arrow[dr, "\mathbb{Z}/2\mathbb{Z}"] & \\
1 \arrow[r, "\mathbb{Z}"] & \mathbb{Z}\omega^2 \arrow[ur, "\Pi_{03}"] \arrow[dr, "\mathbb{Z}/2\mathbb{Z}"'] & & \widetilde{\Pi}_{03} \arrow[ur, "\mathbb{D}_3"] \arrow[dr, "\mathbb{Z}/2\mathbb{Z}"'] & & \widetilde{\mathrm{GL}}(2,\mathbb{Z}) \\
 & & \mathbb{Z}\omega \arrow[ur, "\Pi_{03}"'] & & \widetilde{\Pi}{}^{\tau}_{0,3} \arrow[ur, "\mathbb{D}_3"'] &
\end{tikzcd}

\note{au-dessus du diagramme, trois accolades : $\mathbb{D}'_3$ au-dessus du
segment de $\widetilde{\Pi}{}^{\circ}_{0,3}$ à $\widetilde{\mathrm{SL}}(2,\mathbb{Z})$,
$\mathrm{SL}(2,\mathbb{Z})$ au-dessus du segment de $\mathbb{Z}\omega^2$ à
$\widetilde{\mathrm{SL}}(2,\mathbb{Z})$, $\mathrm{GL}(2,\mathbb{Z})$ au-dessus
du segment de $\mathbb{Z}\omega^2$ à $\widetilde{\mathrm{GL}}(2,\mathbb{Z})$.
Les étiquettes des flèches sont les quotients successifs. La flèche
$\mathbb{Z}\omega \to \widetilde{\Pi}_{03}$ est repassée en gras}

\[
\mathbb{D}'_3 = \{\rho, \sigma \mid \rho^3 = \sigma^2,\ \sigma^4 = 1\ (= \rho^6),\
\sigma(\rho) = \rho^{-1}\}
\]
\[
1 \to \mathbb{Z}/2\mathbb{Z} \to \mathbb{D}'_3 \to \mathbb{D}_3 \to 1
\]

\marginal{NB cette extension n'est pas scindée (les $\rho_{\mathbb{D}'_3}$ qui
relèvent $\rho_{\mathbb{D}_3}$ sont d'ordre 6, \uncertain{non} 3) contrairement
à $\mathbb{D}_6 \to \mathbb{D}_3$ :}

\note{la note est écrite en petit à gauche des suites exactes ; « qui relèvent
$\rho_{\mathbb{D}_3}$ » est inséré par une flèche. Le $\mathbb{Z}$ de
$\mathbb{Z}/2\mathbb{Z}$ est récrit sur une autre lettre}

\begin{tikzcd}[column sep=small, row sep=small, nodes={font=\scriptsize}]
1 \arrow[r] & \mathbb{Z}/2\mathbb{Z} \times \Pi_{0,3} \arrow[r] \arrow[d, no head, "\wr" description] & \mathrm{SL}(2,\mathbb{Z}) \arrow[r] \arrow[d, hook] & \mathbb{D}_3 \arrow[r] \arrow[d, hook] & 1 \\
1 \arrow[r] & \mathbb{Z}/2\mathbb{Z} \times \Pi_{0,3} \arrow[r] & \mathrm{GL}(2,\mathbb{Z}) \arrow[r] & \mathbb{D}_3 \times \{1, \tau\} \arrow[r] & 1
\end{tikzcd}

\note{le $\mathrm{S}$ de $\mathrm{SL}(2,\mathbb{Z})$ est récrit sur un
$\mathrm{G}$}

\begin{tikzcd}[column sep=small, row sep=small, nodes={font=\scriptsize}]
 & & \mathbb{Z}/2\mathbb{Z} \arrow[d] & & \\
1 \arrow[r] & \Pi_{0,3} \arrow[r] \arrow[d, no head, "=" description] & \mathrm{GL}(2,\mathbb{Z}) \arrow[r] \arrow[d, "2"] & \mathbb{D}'_3 \times \{1, \tau\} \arrow[r] \arrow[d, "2"] & 1 \\
1 \arrow[r] & \Pi_{0,3} \arrow[r] & \mathrm{GL}(2,\mathbb{Z})' \arrow[r] \arrow[d, no head, "\wr" description] & \mathbb{D}_3 \times \{1, \tau\} \arrow[r] & 1 \\
 & & \pi_1(\uncertain{U}_{0,3}, \mathbb{D}_3 \times \{1, \tau\}; P_+) & &
\end{tikzcd}

\note{le $\mathrm{G}$ de $\mathrm{GL}(2,\mathbb{Z})'$ est récrit sur une autre
lettre ; dans la dernière ligne, $\mathbb{D}_3$ et $\{1, \tau\}$ sont écrits
par-dessus des ratures. Les deux flèches verticales marquées « $2$ » sont de
degré $2$}

\page{9}

\note{page de schémas : les groupes sont disposés sur la feuille et reliés
par des flèches souvent longues, marquées « cont » (contenu) ou d'un indice ;
on en donne ci-dessous les nœuds et les flèches lisibles, sans reproduire la
disposition. Aucun numéro de sa main}

\[
\Pi^{\mathbb{D}}_{0,3} \longrightarrow \Pi^{\mathbb{D}}_{03}/\Pi'_{0,3}
\xrightarrow{\ \sim\ } \mathbb{Z}/6\mathbb{Z} \times_{\mathbb{Z}/2\mathbb{Z}} \mathbb{D}_3
\simeq (\mathbb{F}_3 \times \mathbb{F}_3) \rtimes \{1, \sigma\}
\longrightarrow \mathbb{Z}/6\mathbb{Z} \longrightarrow \mathbb{Z}/2\mathbb{Z}
\]

\note{sous $\Pi'_{0,3}$, souligné : « sous-groupe d'indice 3 de $\Pi_{0,3}$
obtenu en faisant $\rho_0 = \rho_1 = \rho_\infty$, $\rho_0^3 = 1$ ». Une
accolade marquée $18$ couvre le produit fibré. Le signe devant $\{1, \sigma\}$
est surchargé, lu comme un produit semi-direct}

Au-dessous : $\mathrm{SL}(2,\mathbb{Z}) \to \mathrm{SL}(2,\mathbb{Z})/(\pm 1)
\simeq \mathcal{E}$ (d'ordre $36$) $\to \mathbb{Z}/12\mathbb{Z} \to
\mathbb{Z}/4\mathbb{Z}$, et
$\mathcal{E} \simeq \mathrm{SL}'(2, \mathbb{Z}/6\mathbb{Z})$ (sous-groupe
d'ordre $\uncertain{2}$ à déterminer) ; $\widetilde{\mathrm{SL}}(2,\mathbb{Z})
\to \mathbb{Z}$ ; $\mathbb{D}_3$ (d'ordre $6$) et $\mathbb{D}'_3$ (d'ordre
$12$), avec $\mathbb{D}'_3 \to \mathbb{D}_3$.

\note{le premier terme de la ligne est récrit par-dessus une rature ; lecture
de « $\mathrm{SL}(2,\mathbb{Z})/(\pm 1)$ » incertaine. Les flèches marquées
« cont » vont de $\mathrm{SL}(2,\mathbb{Z})$ vers $\Pi^{\mathbb{D}}_{0,3}$,
de $\mathcal{E}$ vers le produit fibré et vers $\mathbb{D}_3$, de
$\mathbb{Z}/12\mathbb{Z}$ vers $\mathbb{Z}/6\mathbb{Z}$ ; les ordres $6$, $12$,
$18$, $36$ sont écrits près des groupes. Le « $\mathbb{Z}/12\mathbb{Z}$ » au
bout de la flèche partant de $\mathcal{E}$ est écrit sur un
« $\mathbb{Z}/12$ » raturé}

\[
\Pi^{\mathbb{D}}_{0,3} = \{\rho, \sigma \mid \rho^3 = \sigma^2 = 1\}
\qquad
\mathrm{SL}(2,\mathbb{Z}) = \{\rho, \sigma \mid
\underbrace{\rho^3 = \sigma^2}_{\omega},\ \rho^6 = 1\ (= \sigma^4)\}
\]

\[
\Pi^{\mathbb{D}}_{0,3}/\Pi'_{0,3} \simeq \{\rho, \sigma \mid
\rho^3 = \sigma^2 = (\sigma\rho)^6 = 1,\ (\sigma\rho)^2 = (\rho\sigma)^2
\ \struck{\text{central}}\}
\qquad 18
\]

\note{« $= (\rho\sigma)^2$ » est écrit au-dessus du mot biffé, lu
\uncertain{central}}

\[
\Pi^{\mathbb{D}}_{0,3}/\Pi'_{0,3} \xrightarrow{\substack{\rho \mapsto \rho \\ \sigma \mapsto \sigma}}
\mathbb{D}_3 = \{\rho, \sigma \mid \rho^3 = \sigma^2 = 1,\ \sigma(\rho) = \rho^{-1}\}
\quad \text{i.e. } (\sigma\rho)^2 = 1
\]

\[
\Pi^{\mathbb{D}}_{0,3}/\Pi'_{0,3} \longrightarrow \mathbb{Z}/6\mathbb{Z}
\qquad \rho \longmapsto 4 \bmod 6, \quad \sigma \longmapsto 3 \bmod 6
\]

\[
\mathcal{E} = \{\rho, \sigma \mid \underbrace{\rho^3 = \sigma^2}_{\omega},\
\rho^6 = 1\ (= \sigma^4),\ \underbrace{(\sigma\rho)^2\omega}_{\rho_0 = \sigma\rho\sigma^{-1}\rho}
\ \text{invariant par } \rho \text{ et d'ordre } 3\}
\qquad 36
\]
\[
\text{i.e. } (\sigma\rho)^6 = \omega, \quad (\sigma\rho)^2 = (\rho\sigma)^2
\]

\[
\mathcal{E} \longrightarrow \mathbb{Z}/12\mathbb{Z} \simeq (\mathrm{SL}(2,\mathbb{Z}))_{\mathrm{ab}}
\qquad
\begin{aligned}
\rho &\longmapsto 10 \bmod 12 & \sigma\rho = \varepsilon_0 &\longmapsto 1 \bmod 12 \\
\sigma &\longmapsto 3 \bmod 12 & (\sigma\rho)^2\omega = \rho_0 &\longmapsto 8 \bmod 12 \\
\omega &\longmapsto 6 \bmod 12 & &
\end{aligned}
\]

\note{le $10$ est récrit sur un autre chiffre ; entre les deux dernières
lignes de droite, une ligne biffée qui reprend $(\sigma\rho)^2\omega = \rho_0$}

\[
\mathbb{D}'_3 = \{\rho, \sigma \mid \underbrace{\rho^3 = \sigma^2}_{\omega},\
\rho^6 = 1\ (= \sigma^4)\}\ \struck{\ill{}}
\qquad \rho_0 \longmapsto (2 \bmod 6, 1)
\qquad \rho_0 = (\sigma\rho)^2\omega
\]

\note{au-dessus de la relation, un ajout lu « $= (\sigma\rho)^2$ »
\uncertain{}, relié par un trait. Le bas de la page est un second schéma, en
losange : $\mathbb{Z}/6\mathbb{Z} \times_{\mathbb{Z}/2\mathbb{Z}} \mathbb{D}_3$
au sommet, avec au-dessous $\mathcal{E}$ puis $\mathbb{Z}/2\mathbb{Z}$ et
$0$ ; à droite $\mathbb{Z}/6\mathbb{Z}$ (flèche « pr$_1$ ») au-dessus de
$\mathbb{Z}/12\mathbb{Z}$ (flèche marquée « can »), et, plus haut,
$\mathbb{Z}/3\mathbb{Z} \times \{0\}$ $\simeq \mathbb{Z}/3\mathbb{Z}$ ; à
gauche $\mathbb{D}_3$ au-dessus de $\mathbb{D}'_3$. De longues diagonales
relient $\mathbb{D}_3$, $\mathbb{D}'_3$ et $0$ au produit fibré et à
$\mathcal{E}$, et une accolade marquée $(\mathbb{F}_3 \times \mathbb{F}_3)$ et
$\sigma$ court le long de la diagonale principale}

\page{11}\note{p. ④ de l'auteur}

\note{la moitié gauche de la page est écrite en diagonale, à 45 degrés ;
elle est lue ici après rotation. Le texte en est rapide et une bonne part des
mots de liaison ne se lit pas}

\marginal{($n$ impair)}
\textbf{Prop.} \struck{Il} Il existe des relèvements $\tilde\rho$,
$\tilde\sigma$ de $\rho$, $\sigma$ tels que l'on ait
$\tilde\rho^{\,n} = \tilde\sigma^{2}$, et les deux \ill{} sont \ill{}
\ill{} \ill{} $\tilde\omega^{\nu}$ \ill{} l'un des \ill{}
\struck{\ill{} \ill{} \ill{}} \ill{} mod $2n$ (\ill{}) \ill{} \ill{} \ill{}
des choix des relèvements $\tilde\rho$, $(\tilde\sigma)$ et
$\nu \equiv \dfrac{n^2+1}{2}$ \struck{\ill{}}, \ill{}
$\nu \bmod 2n \in \bigl\{\bigl(-\tfrac{n-1}{2}, \tfrac{n+1}{2}\bigr)\bigr\}$
et $\nu$ \uncertain{est impair} ; $\nu_0$ unique pour \ill{} ce qui implique
qu'il existe un \ill{} $\bigl\{-\tfrac{n-1}{2}, \tfrac{n+1}{2}\bigr\}$ qui
soit $\equiv \nu$ $(2n)$. \ill{} $\tilde\rho$, $\tilde\sigma$ sont uniquement
déterminés \ill{} près :
\[
\widetilde{\Pi}{}^{D}_n = \{\tilde\rho, \tilde\sigma, \tilde\omega \mid
\tilde\omega \text{ central},\ \tilde\rho^{\,n} = \tilde\sigma^{2} = \tilde\omega^{\nu}\}
\]

\note{au-dessus du mot « relèvements » de la deuxième ligne, une rature. La
ligne commençant par « \ill{} mod $2n$ » est soulignée et en partie biffée}

À droite, le long de la diagonale :
\[
\mathbb{Z}\tilde\omega^{2} \longrightarrow \widetilde{\Pi}{}^{D}_n
\longrightarrow \Pi^{D}_n = \{\rho, \sigma \mid \rho^n = \sigma^2 = 1\}
\longrightarrow 1
\]

\begin{tikzcd}[column sep=small, row sep=small, nodes={font=\scriptsize}]
 & \mathbb{Z} \arrow[d, "2n"] \\
\Pi^{D}_n \arrow[r] \arrow[dr] & \mathbb{Z} \arrow[d] \\
 & (\Pi^{D}_n)_{\mathrm{ab}} \simeq \mathbb{Z}/2n\mathbb{Z}
\end{tikzcd}

\note{la flèche horizontale part du $\widetilde{\Pi}{}^{D}_n$ de la colonne
diagonale ; la flèche vers $(\Pi^D_n)_{\mathrm{ab}}$ part de
$\Pi^D_n = \{\rho, \sigma \mid \rho^n = \sigma^2 = 1\}$ et est tracée
double. En haut à droite : $\tilde\omega = (1, 2n)$}

\[
\rho \longmapsto (n+1 \bmod 2n) = (1 - n \bmod 2n)
\qquad
\sigma \longmapsto n \bmod 2n
\]

\[
\tilde\rho_0 = (\rho, n+1) \qquad \tilde\sigma_0 = (\sigma, n)
\]
\[
\tilde\rho_0^{\,n} = (\rho^n, n(n+1)) = \struck{\tilde\omega^{\ill{}}}
\qquad
\tilde\sigma_0^{\,2} = (\underbrace{\sigma^2}_{1}, 2n)
\qquad
\tilde\rho_0^{\,n}/\tilde\sigma_0^{\,2} = (1, n(n-1)) = \tilde\omega^{\alpha}
\]
\[
\alpha = \frac{n-1}{2} \qquad (\ill{}\ \alpha + 1)
\]

\[
(\tilde\rho_0\tilde\omega^{\beta})^{n} = (\tilde\sigma_0\tilde\omega^{\gamma})^{2}
\quad \text{i.e.} \quad
\tilde\rho_0^{\,n}\tilde\sigma_0^{-2} = \tilde\omega^{\struck{\ill{}}\,2\gamma - n\beta}
\]

\[
\gamma - n\beta = \alpha = \frac{n-1}{2}
\qquad \struck{\gamma - 2n\beta =}
\qquad \gamma =
\]

\[
(\alpha+1)\cdot 2 - 1 \cdot n = 1
\qquad
\underbrace{\alpha(\alpha+1)\cdot 2}_{\gamma} - \underbrace{\alpha}_{\beta}\cdot n = \alpha
\]

\[
\tilde\rho_0\tilde\omega^{\beta} = \Bigl(\rho,\ n + 1 + 2n\overbrace{\tfrac{n-1}{2}}^{\beta = \alpha}\Bigr)
= (\rho, n^2 + 1)
\]
\[
\tilde\sigma_0\tilde\omega^{\gamma} = \bigl(\sigma,\ n + 2n\underbrace{\gamma}_{\frac{(n-1)(n+1)}{4} = \frac{n^2-1}{4}}\bigr)
= \Bigl(\sigma, \frac{n^3+n}{2}\Bigr)
\]

\note{les deux premiers membres sont écrits avec un $\tilde\omega$ en tête et
un exposant surchargé ($\alpha$, $\gamma$) ; la lecture
$\tilde\rho_0\tilde\omega^{\beta}$, $\tilde\sigma_0\tilde\omega^{\gamma}$ est
celle que demande la suite. Sous la seconde ligne, des accolades
successives donnent $\frac{n^3-n}{2}$ puis $\frac{n^3+n}{2}$ ; le dernier
terme est récrit sur une rature}

\[
\tilde\rho_0^{\,n} = (1, n^3 + n) = \tilde\omega^{\overbrace{\scriptstyle\frac{n^2+1}{2}}^{x}}
\qquad
\tilde\sigma_0^{\,2} = (1, n^3 + n) = \tilde\omega^{\frac{n^2+1}{2}}
\]

\note{dans ces deux lignes, les $\tilde\rho_0$, $\tilde\sigma_0$ sont ceux de
la ligne précédente, corrigés ; le second « $n^3 + n$ » se lit mal}

Séparé par un trait vertical :
\[
\tilde\rho = \tilde\rho_0\,\tilde\omega^{2x} \qquad
\tilde\sigma = \tilde\sigma_0\,\tilde\omega^{nx} \qquad x =
\]

\page{13}\note{p. « 4 bis » de l'auteur}

\[
\tilde\rho^{\,*} = (\rho, n^2 + 1) \qquad
\tilde\sigma^{*} = \Bigl(\sigma, \frac{n^3+n}{2}\Bigr)
\qquad
\tilde\rho^{\,n} \simeq \tilde\omega^{\,?} = \tilde\omega^{\frac{n^2+1}{2}}
\]

\note{l'exposant de $\tilde\rho$ et de $\tilde\sigma$ est un astérisque
appuyé ; le $2$ de $n^2$ est récrit sur un $3$}

\[
(\tilde\rho\,\tilde\omega^{\beta})^n = (\tilde\sigma\,\tilde\omega^{\gamma})^2
\quad \Leftrightarrow \quad
\tilde\omega^{\nu}\,\tilde\omega^{\beta n} = \tilde\omega^{\nu}\,\tilde\omega^{\gamma 2}
\]

\[
\beta n = 2\gamma \qquad
\mathbb{Z}^2 \longrightarrow \mathbb{Z}, \ (\beta, \gamma)
\qquad
\beta = 2x,\ \gamma = nx \qquad \beta n = 2\gamma = (2n)x
\]

$\nu + \underline{2nx}$

\note{un « $\beta$ » biffé devant $\mathbb{Z}^2 \to \mathbb{Z}$. À gauche,
isolés : « $n^2 - n$ » et « $\frac{n-1}{2}$ »}

\[
\frac{n^2+1}{2} \bmod 2n \qquad \frac{n^2+1}{2} \bmod 2n
\qquad \struck{n^2}\ 5 \bmod 6 =
\qquad 13 \bmod 10 \,/\, 3
\]

\[
n = 2\alpha + 1 \qquad
\frac{n^2+1}{2} = \frac{4\alpha^2 + 4\alpha + 2}{2} = \underbrace{2\alpha^2 + 2\alpha + 1}
\ \bmod 4\alpha + 2
\]

\begin{align*}
&(4\alpha+2)\frac{\alpha}{2} + \alpha + 1 && \text{si } \alpha \text{ pair} \\
&\underbrace{(4\alpha+2)\,\frac{\alpha+1}{2}}_{2\alpha^2 + 3\alpha + 1} - \alpha && \text{si } \alpha \text{ impair}
\end{align*}

\[
\frac{n^2+1}{2} \equiv \alpha + 1 = \frac{n+1}{2} \quad (2n) \quad \text{si } \alpha \text{ pair}
\qquad
\frac{n^2+1}{2} \equiv -\alpha = \frac{1-n}{2} \quad (2n) \quad \text{si } \alpha \text{ impair}
\]

\note{le signe $-$ devant $\alpha$ est écrit sur une tache ; le $1$ de
$\frac{1-n}{2}$ est récrit}

\[
\tilde\rho^{\,n} = \tilde\sigma^{2} = \tilde\omega^{\frac{n+1}{2}}
\ \text{si } \frac{n+1}{2} \text{ est impair i.e. } \alpha = \frac{n-1}{2} \text{ pair}
\]
\[
\tilde\rho^{\,n} = \tilde\sigma^{2} = \tilde\omega^{-\frac{n-1}{2}}
\ \text{si } \frac{n-1}{2} \text{ est impair}
\]

donc en tous cas $\tilde\rho^{\,n} = \tilde\sigma^{2} = \tilde\omega^{\nu}$,
$\nu \in \bigl\{-\tfrac{n-1}{2}, \tfrac{n+1}{2}\bigr\}$, $\nu$ impair
--- ce qui suffit : conditions \ill{}.

\note{le premier membre de la seconde ligne est surchargé ; lu
$\tilde\rho^{\,n}$. Le dernier mot est une tache}

\page{14}

\note{écrit sur le verso d'une pièce administrative dactylographiée, tenue à
l'envers ; seul ce qui est de sa main est transcrit. Deux longs traits obliques
traversent la page sans rien biffer}

\[
\Pi^{D}_{03} \longrightarrow
\underbrace{\mathbb{Z}/6\mathbb{Z} \times_{\mathbb{Z}/2\mathbb{Z}} \mathbb{D}_3}_{18}
= \{\nu, g \in \mathbb{Z}/6\mathbb{Z} \times \mathbb{D}_3 \mid \mathrm{sg}^{+}(g) \equiv \nu\ (2)\}
\]

\note{le $\mathbb{Z}/6\mathbb{Z}$ de l'accolade est récrit sur un
$\mathbb{Z}/2\mathbb{Z}$ ; lecture de « $\mathrm{sg}^{+}$ » incertaine}

\struck{et} \quad
$(\Pi_{0,3})_{\mathrm{ab}} \simeq \mathbb{Z}^3/\mathbb{Z}$

\[
(\Pi_{0,3})_{\mathrm{ab}} \otimes_{\mathbb{Z}} \mathbb{F}_3
= \mathbb{F}_3^3/\mathbb{F}_3 \supset \ill{}(\mathbb{F}_3^{3\,\uncertain{*}})/\mathbb{F}_3
\]

\[
\rho_0\ \rho_1\ \rho_\infty \quad
\begin{cases}
\rho_0\rho_1^{-1},\ \rho_1\rho_\infty^{-1},\ \rho_\infty\rho_0^{-1},\ \rho_0^3 \\
\struck{[\rho_0, \rho_1]\,[\rho_1, \rho_\infty]\,[\rho_\infty, \rho_0]}
\end{cases}
\]

\[
\mathrm{SL}(2, \mathbb{Z}/6\mathbb{Z})/\{\pm 1\}
\simeq \bigl(\underbrace{\mathrm{SL}(2, \struck{\ill{}}\,\mathbb{Z}/2\mathbb{Z})}_{6} \times
\underbrace{\mathrm{SL}(2, \mathbb{Z}/3\mathbb{Z})}_{24}\bigr)/\pm 1
\qquad \text{d'ordre } 72
\]

\note{sous le $6$, une petite ligne lue \uncertain{$p(p-1)(p+1)$} ; le
« d'ordre » est souligné, et l'accolade couvre les deux membres}

\[
\{\rho, \sigma \mid \rho^3 = \sigma^2 = \struck{\ill{}}\,(\sigma\rho)^6 = 1,\
\sigma\rho\sigma\rho = \rho\sigma\rho\sigma\}
\]

\note{au-dessus de la dernière relation : « $(\sigma\rho)^2 = \rho((\sigma\rho)^2)$ »}

\[
\rho_0 = \sigma\rho\sigma\rho \qquad \rho_1 = \rho\sigma\rho\sigma
\qquad
\struck{[\rho_0, \rho_1] = \ldots = [\rho_1, \rho_0]}
\]

\note{le commutateur biffé est développé sur deux lignes illisibles sous la
rature ; au-dessous, un produit isolé lu
\uncertain{$\sigma\rho\sigma\rho\sigma^{-1}\rho^{-1}\sigma^{-1}\rho^{-1}$}, et
en marge « $\rho_\infty = \uncertain{\sigma\rho}$ »}

Au pied de la page :
\[
\rho_0 \quad \rho_1 \quad \rho_\infty \quad 5
\qquad
\rho_0 = \struck{\ill{}}\,\sigma\rho\sigma^{-1}\rho = (\sigma\rho)^2\omega
\]

\note{en marge de gauche, écrits de biais : $\sigma\rho\sigma\rho = 1$ et
une colonne de puissances lues « $\uncertain{\sigma\rho^{-1}\ 3}$,
$\uncertain{9\rho\ 4}$, $\uncertain{9\rho\ 2}$, $2\cdot 4$ », reliées par une
flèche marquée $\rho$}

\page{15}\note{p. 5 de l'auteur}

\marginal{à vérifier --- OK}
$M_{1,1} \simeq (U_{0,3}, \mathbb{D}'_3)$, où $\mathbb{D}'_3$ est défini
comme produit fibré

\begin{tikzcd}
\mathbb{D}_3 = \mathfrak{S}_3 \arrow[r, "\mathrm{can}"] & \mathbb{Z}/2\mathbb{Z} \\
\mathbb{D}'_3 \arrow[u] \arrow[r] & \mathbb{Z}/4\mathbb{Z} \arrow[u, "\mathrm{can}"']
\end{tikzcd}

\note{la note marginale, à gauche, est écrite de biais ; « à vérifier » est
barré de deux traits, « OK » écrit au-dessous. Le $U$ de $U_{0,3}$ est une
capitale ronde, lue $U$ comme au folio 7}

donc $\mathbb{D}'_3$ est une extension (\uncertain{bien} \uncertain{définie})
de $\mathbb{D}_3$ par $\mathbb{Z}/2\mathbb{Z}$ ; \struck{\uncertain{beaucoup}}
on fait opérer $\mathbb{D}'_3$ sur $U_{0,3}$ via $\mathbb{D}_3$, donc avec
\uncertain{trivialité} de l'élément $\omega$ d'ordre $2$ du noyau de
$\mathbb{D}'_3 \to \mathbb{D}_3$.

\note{le dernier groupe est écrit « $\mathbb{D}'_2 \to \mathbb{D}_2$ », les
indices mal formés ; le contexte demande $\mathbb{D}'_3 \to \mathbb{D}_3$}

Or a donc
\[
\mathcal{T}_{11} \simeq \pi_1(\overbrace{U_{0,3}, \mathbb{D}'_3}^{\Pi_{0,3}^{\mathbb{D}'_3}}; P_+)
\quad \text{et ainsi}
\qquad
\mathcal{T}_{11} \simeq \mathrm{SL}(2,\mathbb{Z})
\]

\note{« $\simeq \mathrm{SL}(2,\mathbb{Z})$ » est écrit sous $\mathcal{T}_{11}$,
précédé d'un signe $\wr$ ; de l'isomorphisme part un trait vers une note :
« dépend du choix d'un pt du demi-plan de Poincaré au-dessus de $P_+$, on peut
prendre $\zeta = \exp(2i\pi/6)$ ». À gauche de la note, un croquis : le cercle
unité dans le plan, un point $\zeta$ sur le cercle, relié à l'origine, avec
l'abscisse $\frac12$ marquée}

\begin{tikzcd}
\Pi^{\mathbb{D}'}_{0,3} \arrow[r] \arrow[d] & \Pi^{\mathbb{D}}_{0,3} = \{\rho, \sigma \mid \rho^3 = \sigma^2 = 1\} \arrow[d, Rightarrow] \\
\mathbb{D}'_3 \arrow[r] \arrow[d] & \mathbb{D}_3 \arrow[d] \\
\mathbb{Z}/4\mathbb{Z} \arrow[r, no head] & \mathbb{Z}/2\mathbb{Z}
\end{tikzcd}

\note{les deux carrés sont marqués « cart » (cartésiens). Au-dessus de
$\Pi^{\mathbb{D}'}_{0,3}$, « $\Pi_0 \subset \mathbb{D}'$ » \uncertain{} ;
$\mathbb{D}'_3$ et $\mathbb{D}_3$ sont écrits par-dessus d'autres symboles
(« $\mathbb{Z}/4$ », « $\mathbb{Z}/2\mathbb{Z}$ ») raturés}

À droite du diagramme :
\struck{$\rho \mapsto 0$, $\sigma \mapsto 1$} \quad $\rho \mapsto 0$,
$\sigma \mapsto 1$

On remplace $\rho$, $\sigma$ par $\rho'$, $\sigma'$ :
\[
\rho' = (\rho, 2), \quad \sigma' = (\sigma, 3) \qquad
\text{alors } \rho'^{6} = \sigma'^{4} = 1, \quad
\rho'^{3} = \sigma'^{2} = (\underbrace{1, 2 \bmod 4}_{\omega})
\]

\note{les secondes composantes $2$ et $3$ sont récrites sur d'autres chiffres}

\marginal{Bien vérifié ainsi, notamment que
$\mathrm{SL}(2,\mathbb{Z}) \simeq \{\rho', \sigma' \mid \rho'^{3} = \sigma'^{2},
\ \ldots\}$}

\note{note écrite de biais dans la marge de droite, reliée par une longue
accolade à l'accolade de $\rho'$, $\sigma'$ ; la fin des relations se lit mal}

Ainsi $\Pi^{\mathbb{D}'}_{0,3} \simeq \mathrm{SL}(2,\mathbb{Z})
\ \uncertain{apparaît} comme extension de $\Pi^{\mathbb{D}}_{0,3}$ par
$\mathbb{Z}/2\mathbb{Z}$, \uncertain{image} \uncertain{inverse} de l'ext.
$\mathbb{D}'_3$ de $\mathbb{D}_3$ par $\mathbb{Z}/2\mathbb{Z}$, \ill{} de
l'ext. $\mathbb{Z}/4\mathbb{Z}$ de $\mathbb{Z}/2\mathbb{Z}$ par
$\mathbb{Z}/2\mathbb{Z}$, via $\Pi^{\mathbb{D}}_{0,3} \to \mathbb{D}_3$
\ill{} $\Pi^{\mathbb{D}}_{0,3} \to \struck{\mathbb{D}_3}\, \mathbb{Z}/2\mathbb{Z}$.
C'est une ext. de $\mathbb{D}'_3$ par un \uncertain{sous-groupe}
$\simeq \Pi_{0,3}$, \uncertain{choisi} $\rho_0$, $\rho_1$, $\rho_\infty$ \ldots

\note{la page s'arrête sur ces points de suspension, au ras du bord}

\section{Sommets, arêtes, faces}

\note{titre de l'édition. Page 16 à l'encre bleue ; un long trait de crayon la
traverse en diagonale sans rien biffer}

\page{16}

\[
\vec{S}(X_\bullet) \simeq G_0 \qquad
S(X_\bullet) \simeq G_0/H_S, \quad A(X_\bullet) \simeq G_0/H_A, \quad
F(X_\bullet) \simeq G_0/H_F
\]

\note{de $\vec{S}(X_\bullet) \simeq G_0$ partent trois flèches, vers chacun des
trois quotients}

\[
H_S \cap H_A = e, \quad H_A \cap H_F = e, \quad H_S \cap H_F = e
\]
i.e.
\[
G_0 \hookrightarrow G_0/H_S \times G_0/H_A, \quad
G_0 \hookrightarrow G_0/H_A \times G_0/H_F, \quad
G_0 \hookrightarrow \struck{\ill{}}\,G_0/H_A \times G_0/H_F
\]

\note{la troisième injection répète la deuxième ; on attendrait
$G_0/H_S \times G_0/H_F$. Son premier facteur est récrit sur une rature}

\marginal{Cas non orienté}
Considérons la catégorie des polyèdres (non orientés/des \ill{}) \ill{}
type donné, il y a des foncteurs
\[
S \quad \vec{A} \quad F \quad \struck{\ill{}}
\]

\note{« $\vec{A}$ » : un $A$ surmonté d'une flèche, sous un $S$ ; lecture
incertaine}

\begin{itemize}
\item[$A_1$] relation d'incidence entre $S$ et $A$ (arêtes orientées au
sens 1)
\item[$A_2$] --- --- $A$ et $F$ (--- --- 2)
\item[$A_3$] rel. d'incidence entre $S$ et $F$
\end{itemize}

$R = \struck{\ill{}}\ A_1 \times_A A_2$ \quad
\struck{rel. d'incidence entre $S$ et $F$} drapeaux des
$(s, a, f)$, $s < a < f$.

\note{« drapeaux des » est écrit au-dessus de la ligne biffée}

\begin{tikzcd}
 & R \arrow[dl] \arrow[d] \arrow[dr] & \\
S & A & F
\end{tikzcd}

\[
A_1 = \mathrm{Im}(R \to S \times A) \qquad
A_2 = \mathrm{Im}(R \to A \times F) \qquad
A_3 = \mathrm{Im}(R \to S \times F)
\]

$\Omega$ ens. des orientations

\[
A_{1\,\Omega} \simeq A_{2\,\Omega} \simeq A_{3\,\Omega} \quad (
\]

\note{la parenthèse reste ouverte. Le $A_2$ est récrit sur un $A_3$}

\[
\mathrm{Pol}_\tau \qquad \mathrm{Pol}_\tau^{\Omega} \simeq \mathrm{Pol}_{\tau\,/\Omega}
\]

\note{l'exposant de $\mathrm{Pol}_\tau$ est récrit sur une rature, lu
$\Omega$ ; l'indice du dernier terme est écrit sur une rature}

\section{Le groupe $\mathcal{T}_{1,1}$}

\note{titre de l'édition ; la suite numérotée reprend (pages 17 et 19)}

\page{17}\note{p. ⑥ de l'auteur}

\[
S\mathcal{T}_{1,1} \simeq \text{image inverse de } \mathrm{SL}(2,\mathbb{Z})
\text{ dans } \widetilde{\mathrm{SL}(2,\mathbb{R})}
\simeq \{\rho, \sigma \mid \rho^3 = \sigma^2\ (= \omega)\}
\qquad (\omega \text{ élément central})
\]

\[
\mathcal{T}_{11} = S\mathcal{T}_{11}/\langle\omega^2\rangle
\]

\note{le $\mathrm{SL}(2,\mathbb{R})$ est écrit sur une rature, le tilde à la
fin. Le quotient est noté « $S\mathcal{T}_{11}\ k\omega^2\rangle$ », le
chevron ouvrant mal formé}

\[
(\uncertain{S}\mathcal{T}_{11})_{\mathrm{ab}} \simeq \mathbb{Z}
\qquad
\struck{(\rho, \sigma)}\quad
\rho \longmapsto 2, \quad \sigma \longmapsto 3, \quad \omega \longmapsto 6
\]
\[
(\mathcal{T}_{11})_{\mathrm{ab}} \simeq \mathbb{Z}/12\mathbb{Z}
\qquad
\rho \longmapsto 2 \bmod 12, \quad \sigma \longmapsto 3 \bmod 12
\]

\note{le premier groupe est écrit « $(\mathcal{T}_{11})_{\mathrm{ab}}$ », sans le
$S$ que demandent les valeurs $2$, $3$, $6$ dans $\mathbb{Z}$ ; le $6$ est
récrit sur un autre chiffre}

\begin{tikzcd}
1 & & 1 \\
\mathcal{T}_{11} \arrow[u] \arrow[rr, Rightarrow, "\struck{\ill{}}"] & & \mathbb{Z}/12\mathbb{Z} \arrow[u] \\
S\mathcal{T}_{11} \arrow[u] \arrow[rr] & & \mathbb{Z} \arrow[u] \\
\mathbb{Z} \arrow[u] \arrow[rr, "\mathrm{id}"] & & \mathbb{Z} \arrow[u, "12"'] \\
1 \arrow[u] & & 1 \arrow[u]
\end{tikzcd}

\note{la flèche du haut porte un mot raturé, puis un « \struck{continûment} »
\uncertain{} sous elle. À gauche du $\mathbb{Z}$ de gauche :
« $\varpi = \omega^{2}$ » (exposant surchargé), au-dessus d'une flèche partant
de $1$}

Donc $S\mathcal{T}_{11}$ se déduit de $\mathcal{T}_{11}$ par image inverse de
l'ext. canonique de $\mathbb{Z}/12\mathbb{Z}$ par $\mathbb{Z}$, via
$\mathcal{T}_{11} \to \mathbb{Z}/12\mathbb{Z}$.

\note{le bas de la page porte, tête-bêche et à l'encre bleue, un fragment
d'une autre rédaction, barré de trois longs traits de crayon ; il est donné
ici tel qu'il se lit, feuille retournée}

\struck{\ill{} normalisations de $(K, K_0)$, dans elles est équivalente à la
cat. des variétés dans $K_f N$ --- \ill{} \ill{} revêt\ill{} top. --- une
unique normalisation pour laquelle la projection soit un morphisme de
$\Delta$-complexes normalisés. [Mais attention, \ill{} \ill{} \ill{} pas
\uncertain{structures} normalisées \ill{} \ill{}\ldots]}

\struck{8. Groupes et groupoïdes fondamentaux. On montrera la relation entre
construction top. et construction algébrique.}

%P18%
\end{document}
